\section{Background: kernel methods and random features}

Kernel methods are a broad class of machine learning algorithms
which make predictions using a positive definite \emph{kernel function}
$k:\mathcal{X}\times\mathcal{X}\mapsto\mathbb{R}$
\citep{scholkopf2002learning}.
Common methods in this class are support vector machines \citep{cortes1995support}
and Gaussian processes \citep{williams2006gaussian}.
Given a dataset of $n$ data points,
training most kernel methods requires computing 
the $n\times n$ kernel matrix\footnote{
We write $x_i$ to refer to the $i$th element of a vector and $x^{(i)}$ to denote the $i$th vector in a list.
}
$K_{i,j} = k(x^{(i)},x^{(j)})$ (with $O(n^2)$ time complexity)
and possibly inverting it (with $O(n^3)$ time complexity).
Because of this, applying kernel methods to large datasets generally requires approximations.

Given a kernel $k$, a \emph{random features map} 
is a random function
$f:\mathcal{X}\mapsto\mathbb{R}^M$, with the property that  $f(x)\cdot f(x')$ approximates $k(x,x')$ for every pair $x,x'\in\mathcal X$. The approximation is often exact in expectation:
\begin{equation}\label{eqn:random-feature-definition}
    \mathbb{E}_f\left[f(x)\cdot f(x')\right]=k(x,x') \quad \text{for all }x,x'\in\mathcal X.
\end{equation}
Random features allow the kernel matrix to be approximated as 
$\widehat K_{i,j} = f(x^{(i)})\cdot f(x^{(j})$.
Because this matrix has rank at most $M$,
this approximation generally reduces the cost of $O(n^3)/O(n^2)$ computations to $O(M^3)/O(nM^2)$,
i.e.\@ at most linear in $n$.



The seminal work of \citet{rahimi2007random}, which coined the term random features, gave a general method based on Fourier analysis to construct random features for any \emph{Bochner} or \emph{stationary kernel}, for which $k(x,x')$ is a function of $x-x'$. This class includes many common kernels including the RBF and Mat\'{e}rn kernels, but excludes $T_{MM}$.
Subsequent works have proposed random features for other kernels including the polynomial kernel and the arc-cosine kernel 
\citep{liu2021random}.
However, there is no general formula to define random features for non-stationary kernels, such as $T_{MM}$.

A random features map is sometimes called a \emph{data-oblivious sketch}, to distinguish it from other \emph{data-dependent} low-rank approximation methods which depend on a given dataset $x^{(1)},\dots,x^{(n)}$.
Examples of data-dependent low rank sketches are the Nystr\"{o}m approximation and leverage-score sampling \citep{drineas2005nystrom,drineas2012fast}.
Although data-dependent methods may result in lower approximation errors for a given dataset,
data-oblivious sketches are naturally parallelizable and useful in cases where the dataset  changes over time (e.g.\@ streaming or optimization)
or for ultra-large datasets which may not fit in memory.



